\chapter{模型推断与生成}
\label{chap:nlp-model-inference}

同一个已经完成微调和对齐的模型，可以为同一份通知生成几种不同的参加指南。变化可能来自提示里的要求，也可能来自抽样时的随机选择；即使输入完全相同，最大生成长度不同，也可能使一份指南完整而另一份停在半句。解释这些差异，需要把参数形成过程与输出形成过程分开。前面三章研究怎样学习参数，本章固定参数$\bm{\theta}$，考察一次调用怎样把条件序列变成可交付结果。

这里的“推断”对应inference，指使用模型计算并形成输出；下一章讨论的“推理”对应reasoning，指依据条件完成多步求解。一次普通翻译也需要模型推断，推理任务同样要经过本章的输入组织、状态计算、选择和终止环节。增加运行时计算可能改变求解表现，但不会自动改写模型参数。

以下用一份教学通知串联输入和输出：“读书会改至4月18日14:00，在海棠楼201举行；校外来宾须在4月16日18:00前提交报名信息。”要求是为校外来宾生成两句参加指南。除明确引用的论文外，词元分布、候选和预算数值均为教学构造，用来核对机制，不是模型测量。

\section{输入与上下文组织}
\label{sec:nlp-inference-context}

模型读取的是词元序列，而应用接到的是角色消息、用户要求和若干材料。首先要确定这些对象怎样进入同一个条件，使后续概率计算确实针对当前任务。输入组织会改变模型看到的证据及要求，因此本身就是生成规则的一部分。

\subsection{指令与对话序列组织}
\label{sec:nlp-inference-dialogue}

假设用户先问活动地点，得到回答后又说“再写上报名截止时间”。第二次请求里的“再”依赖前一轮回答；只保留最后一句，就丢掉了待修改对象。可以沿用第\ref{chap:nlp-supervised-finetuning}章的对话格式，把系统要求、用户消息和历史回答依次序列化，末尾放置待生成回答的起始标记。角色及消息边界必须采用该模型学习过的约定，不能仅凭界面上显示了不同颜色，就假定模型能够区分消息来源。

将最终条件序列记为$c=(c_1,\ldots,c_n)$。其中可以包含“保留精确日期和报名条件”的要求、原始通知、前一轮回答与当前修改请求。模型接到的条件还包括位置编号和注意可见关系；这两者相同，才是在比较相同输入。对话模板若自动补充日期、工具说明或特殊标记，也要计入$c$及其长度。

角色标记提供来源信息，并不构成不可违背的执行权限。模型仍可能把引用材料中的指令当成当前要求。因而组织输入时应明确标出待处理内容的范围，并把任务要求与材料文本分别说明；后文的结构检查也不能替代这种来源区分。

\subsection{示例与材料的条件化组织}
\label{sec:nlp-inference-demonstrations}

已经明确消息来源后，还需要告诉模型当前输出承担什么职责。可以加入一条不同活动的“通知—两句指南”示例，示范应保留哪些字段、怎样表达外部参加条件。示例在本次推断期间进入$c$，参数仍固定；这与使用这些样本执行有监督微调有不同的操作对象。示例不应包含当前测试题的参考答案，否则比较中就混入了答案信息。

外部材料也通过条件序列发挥作用。例如，把报名规则和活动通知分别放在带来源名称的引用区间中，再要求合并生成指南。模型能够据此形成条件概率，但材料被放入上下文并不保证它被正确使用。两份材料对地点有冲突时，单纯增加提示长度不能决定哪个来源可信，必须先有来源优先级或澄清规则。

示例、要求与证据的顺序会改变具体序列$c$。因此比较不同模型时，应记录相同的材料集合如何排列，而不只说“都给了同一份资料”。检索怎样取得材料属于后面的问答应用；此处只处理已经取得的材料怎样成为输入条件。

\subsection{上下文与生成预算分配}
\label{sec:nlp-inference-budget}

材料放入之后，还要为尚未生成的回答预留空间。设一次调用允许的总长度为$L$，输入占$n$个词元，生成上限为$m$，最简单的约束是$n+m\leq L$。角色模板、工具说明和结束标记也消耗位置；只计算可见汉字数，会低估实际词元长度。若一个模型分别限制输入与输出长度，还须同时满足各自的限制。

超过预算时，可以删去无关材料、选择证据片段，或把较早对话压缩成状态摘要。这些操作都会改变条件。若摘要只保留活动时间而删去“校外来宾”的限定，后续就可能把报名要求施加给所有人；这是输入条件发生变化，不是解码阶段能够补回的信息。压缩结果应保留实体、否定、时间和例外的绑定关系。

允许放入$L$个位置，也不意味着模型在所有位置都能同样可靠地利用证据。长度适应的训练依据见第\ref{chap:nlp-large-language-models}章；在当前调用中，还应检查关键材料是否保留，以及生成上限是否足以完成任务。整理后的$c$与预算一起进入概率计算。

\section{概率计算与增量生成}
\label{sec:nlp-inference-computation}

确定输入以后，模型计算给定条件下的词元概率。这里先说明概率从何而来、哪些中间量能够复用；概率最高的词元是否被选中，则由下一节的解码规则决定。

\subsection{上下文计算与逐步回填}
\label{sec:nlp-inference-autoregression}

对因果语言模型，把输入$c$完整送入前向计算，可以得到输入末端的状态及下一词元分布。设已经生成的前缀为$y_{<t}$，本步分布为
\begin{equation}
 p_t(v)=p_{\bm{\theta}}(v\mid c,y_{<t}),\qquad v\in V,
 \label{eq:nlp-inference-next-token}
\end{equation}
其中$V$是词表。某条规则选出$y_t$以后，将它接到前缀末端，再计算$p_{t+1}$。这一循环持续到结束条件触发。训练时参考序列已知，多个监督位置可以一起计算；运行时尚不知道要选择的未来词元，不能把它们当成已知历史。

计算初始上下文通常称为预填充，即把已有序列转换成后续能够读取的状态。逐步回填则每次增加已经获准的输出。两者调用同一套已学习参数，但输入长度和可复用状态不同。概率计算既可以服务于贪心选择，也可以服务于抽样或候选评分，本身没有规定最终生成哪条序列。

\subsection{历史状态的缓存与更新}
\label{sec:nlp-inference-cache}

若每生成一个词元都重新处理整个$c,y_{<t}$，相同前缀会被反复计算。因果注意中的历史位置不能读取后来追加的词元，因此在参数、前缀、位置和可见关系都未改变时，历史各层的键和值仍与上一步相同。可以缓存它们，只计算新位置的查询、键和值。选择$y_t$并把它回填到位置$N=n+t$以后，计算该位置的表示，以便预测$y_{t+1}$。对某一层，单头的计算可写成
\begin{equation}
 \bm{o}_N=\sum_{j=1}^{N}\alpha_{N,j}\bm{v}_j,
 \quad
 \alpha_{N,j}=\frac{\exp(\bm{q}_N^{\top}\bm{k}_j/\sqrt{d_k})}
 {\sum_{r=1}^{N}\exp(\bm{q}_N^{\top}\bm{k}_r/\sqrt{d_k})},
 \label{eq:nlp-inference-cache-attention}
\end{equation}
其中$\bm{q}_N,\bm{k}_j\in\mathbb{R}^{d_k}$，$\bm{v}_j\in\mathbb{R}^{d_v}$是值向量，历史的$\bm{k}_j,\bm{v}_j$从缓存读取。$N$表示回填后的实际序列长度；预测$y_t$用的是回填前的位置$N-1$的表示，因而没有提前读入$y_t$。缓存减少了重复计算，但当前查询仍要读取允许的历史状态。

对每层缓存$H_{\mathrm{kv}}$组键值、每组维数$d_k,d_v$的普通注意，长度为$N$时缓存标量数约为$N H_{\mathrm{kv}}(d_k+d_v)$，再乘层数。第\ref{chap:nlp-large-language-models}章的GQA减少共享键值组的数量，MLA则缓存压缩表示；二者改变缓存对象的参数化，不能都按独立全尺寸键值组估算。递归状态模型缓存的是状态而不是全部历史键值，其可表达的记忆机制也不同。

修改系统要求、插入历史消息、重新编号位置，或者改变上下文窗口，都会使部分缓存失去原先的依据。安全的复用范围是经过核对的共同前缀；若采用近似量化或状态压缩，还须把近似误差单独考虑。精确缓存与完整前向在相同计算约定下对应同一条件分布，有限精度和实现次序可能带来数值差异。

缓存本身没有学习目标，也不增加待拟合参数：键和值由已训练的投影矩阵及当前前缀计算得到。设有$D$层、批量为$b$、每个标量占$s$字节，普通键值缓存约占$bDNH_{\mathrm{kv}}(d_k+d_v)s$字节。每步只形成一个新位置，仍需读取长度随$N$增长的历史；不能把省去历史前馈和投影的代价，说成整步注意力变为常数时间。多个候选共享前缀时可以共享只读缓存，但某一候选追加的状态不能直接接给另一候选。

长短不一的请求还会带来存储分配问题。按最大生成长度提前保留连续空间，会留下尚未使用的槽位；束搜索复制共同前缀，又可能重复存储同一状态。PagedAttention把逻辑词元位置映射到固定大小的物理块，通过块表读取非连续键值，共同前缀可共享块，发生分支写入时再复制。块表和引用计数是运行时记录，不是从样本学出的模型参数。原论文图2在其服务负载中报告，三种Orca配置实际词元状态只占预留键值空间的$20.4\%$至$38.2\%$，vLLM约为$96.3\%$；论文的吞吐收益来自缓存管理、批处理与调度的组合，不是语言准确率增加。它适用于同一模型服务多个长度不同的请求，不能据此推断单条短指南也会按同一比例加速\scite{kwon2023pagedattention}。

\subsection{掩码序列的迭代恢复}
\label{sec:nlp-inference-masked-restoration}

因果模型按输出位置推进，而掩码语言生成模型可以从若干未确定位置开始，反复恢复整段序列。沿用第\ref{chap:nlp-large-language-models}章的掩码去噪模型，先为待生成区域放入掩码，再在当前可见文本条件下预测各掩码位置。每轮按指定日程确定一部分词元，也可以重新掩盖低置信位置，再执行下一轮。迭代次数、每轮确定多少位置，以及是否分块，都参与生成规则。

LLaDA与Dream提供了这类生成路径的实例；其训练目标与具体恢复日程应分别阅读，不能因为都采用掩码就假定推断过程相同。与因果回填相比，它们允许一个位置利用已显露的左右信息；某个词元被替换后，其他位置的状态也可能改变。因此上一小节关于“历史不读未来”的缓存理由不能直接照搬。多位置预测也不等于彼此独立地一次抽取即可获得目标联合分布。

具体运行时，先固定条件$c$与待生成长度$m$，记第$r$轮的待恢复序列为$\tilde y^{(r)}$。模型为其中每个掩码位置$j$给出$p_{\bm\theta,j}(v\mid c,\tilde y^{(r)})$，按采样规则提出临时词元，再按本轮配额决定保留哪些位置。LLaDA的低置信度重掩码方案把低概率提议重新变成掩码，使较有把握的位置先成为下一轮条件；其半自回归方案则依次恢复若干块，块内仍可同时处理多个位置。这里的重掩码不必意味着每轮推翻所有此前已经确定的词元，应以实际日程规定的可修改集合为准。配额若为零，该轮不能推进；最后一轮还须处理剩余掩码，而不是把它们作为正文交付。

这些恢复概率来自第三章的随机损坏与加权去噪训练，不是采样器在本次请求中重新估计的参数。LLaDA从头训练掩码预测器，指令阶段只对回答区域损坏并学习恢复；Dream从已有自回归模型初始化，保留位置$i$的表示预测位置$i+1$的位移约定，再适应全注意力和去噪目标。因而从Dream输出张量取某个位置的分布时，必须按其位移恢复目标位置，不能直接套用LLaDA的同位置接口。轮数、块长、温度和置信度阈值是运行超参数，可在独立验证材料上选择，不具有自己的反向传播步骤。

这种路径的实际用途是从两端已有文字恢复中间内容，或在给定槽位中联合组织输出。LLaDA原论文表1（第6页）报告，其8B基座在两示例HumanEval-FIM任务上的通过率为$73.8\%$；普通零示例HumanEval为$33.5\%$。两项任务的可见信息和提示不同，不能把差值称为同任务加速收益，它们展示的是补全接口可以利用后缀约束。原论文表3的诗句正向、反向补全也提供了双向条件使用的观察。与逐词生成相比，增加去噪轮次可能改善恢复，但每轮通常重算多个位置，较少轮次又可能过早固定不相容的词元；并行位置数不是实际速度的充分依据\scite{nlp-llm-llada2025}。

Dream的实验进一步比较联合规划输出。原论文表1（第7页）在相同八示例Sudoku协议下，Dream 7B基座为$81.0\%$，Qwen2.5 7B基座为$21.0\%$；同表MMLU则为$69.5\%$与$71.9\%$。这些是不同模型完整训练与生成方案的比较，不是仅替换采样日程的消融，也不代表扩散在所有任务上胜出。报告第五节的用途是约束满足式规划，第六节进一步研究生成顺序和去噪轮次；一般指令与代码仍须分别检查\scite{nlp-llm-dream2025}。

\section{解码与采样}
\label{sec:nlp-generation-candidates}

模型分布只是对延续的赋值，应用还要决定从中取什么。搜索希望得到高评分序列，抽样希望按照指定分布保留多种可能；这两个目标不同。输出的长度、格式与完成状态持续约束选择过程，不能等到最后才决定它们。

\subsection{高概率序列的近似搜索}
\label{sec:nlp-generation-beam}

以下以EOS表示结束词元，以BOS表示启动生成时使用的起始标记；它们不作为普通文字交付。\term{贪心解码}（greedy decoding）每一步只选概率最高的后继。
\term{束搜索}（beam search）同时保留$b$个尚未结束的前缀，以降低早期选择失误的影响。
用累计对数概率$\ell(h)$给前缀$h$计分，扩展一个词元$v$时计算
\begin{equation}
  \ell(hv)=\ell(h)+\log p_{\bm{\theta}}(v\mid c,h),
  \qquad s(h)=\frac{\ell(h)}{|h|^\alpha},\quad \alpha\geq0.
  \label{eq:nlp-generation-beam-score}
\end{equation}
$|h|$不含BOS，完整候选的长度含EOS；空前缀只用于初始化，不代入归一化式。
本节用幂次长度归一化作为明确的教学约定，并非所有工具包的默认公式。
同一步的活动前缀等长，归一化不改变排序；不同长度的完成候选则可能换序。
由于对数概率非正，$\alpha>0$会减轻长序列的累计负分，但不能保证恰当长度。

\begin{algorithm}[带独立结束集合的束搜索]
  \label{alg:nlp-generation-beam}
  \begin{algorithmic}[1]
    \Require 条件$c$，束宽$b$，最大生成步数$m$，指数$\alpha$
    \Ensure 完成候选，或无完成候选的失败状态
    \State 活动集合$H\gets\{(\epsilon,0)\}$；结束集合$F\gets\varnothing$
    \For{$t=1,\ldots,m$}
      \State $C\gets\varnothing$
      \For{每个$(h,\ell)\in H$及词表中的每个$v$}
        \State $\ell'\gets\ell+\log p_{\bm{\theta}}(v\mid c,h)$
        \If{$v=\mathrm{EOS}$且满足结束条件}
          \State 把$(hv,\ell')$加入$F$，以后不再扩展
        \ElsIf{$v\ne\mathrm{EOS}$且前缀合法}
          \State 把$(hv,\ell')$加入$C$
        \EndIf
      \EndFor
      \State $H\gets C$中按$s$最高的至多$b$项；平局按预定词元序
      \State $F\gets F$中按$s$最高的至多$b$项
      \If{$H=\varnothing$}
        \State 结束循环
      \EndIf
    \EndFor
    \State 返回$F$中的最佳项；若$F=\varnothing$，报告预算内未完成
  \end{algorithmic}
\end{algorithm}

算法~\ref{alg:nlp-generation-beam}不把结束候选重新放回活动束，也不把达到长度上限的
未结束前缀伪装成完整文本。它最多执行$m$轮，每轮至多扩展$b|V|$个符号，
这里$|V|$是词表大小，模型前向计算的代价还需另计。
为便于核对，算法不采用提前终止。若$\alpha=0$，未来追加的对数概率不会增加累计分数，
当前最佳结束分数不低于所有活动分数时就可提前停止寻找最佳项。
若使用长度归一化，不能继续套用这个判断；有限$m$时，活动前缀的
$\ell(h)/m^\alpha$可作为此归一化下的乐观上界，但保证也仅相对于尚未剪掉的前缀。

\Needspace{12\baselineskip}
\begin{example}[三步生成中的一次早期遗漏]
  \label{ex:nlp-generation-beam}
  词表只有$A,B,\mathrm{EOS}$，$A,B$是抽象文字片段；最多生成三步，EOS占一步。
  下表给出所有会影响计算的条件分布，其他前缀统一采用$(0.10,0.10,0.80)$。
  \begin{center}
    \begin{tabular}{@{}lrrr@{}}
      \toprule
      已有前缀 & $A$ & $B$ & EOS\\
      \midrule
      空 & 0.55 & 0.40 & 0.05\\
      $A$ & 0.45 & 0.35 & 0.20\\
      $B$ & 0.10 & 0.10 & 0.80\\
      $AA$ & 0.35 & 0.25 & 0.40\\
      $AB$ & 0.10 & 0.10 & 0.80\\
      \bottomrule
    \end{tabular}
  \end{center}
  贪心依次选择$A,A,\mathrm{EOS}$，完整概率为
  $0.55\times0.45\times0.40=0.099$。
  取$b=2,\alpha=0$时，第一步活动束为$A:0.55,B:0.40$，
  空文本的EOS候选$0.05$另存。第二步从$A$得到
  $AA:0.2475,AB:0.1925,A\mathrm{EOS}:0.11$；
  从$B$得到$BA:0.04,BB:0.04,B\mathrm{EOS}:0.32$。
  活动束保留$AA,AB$，结束集合保留$B\mathrm{EOS},A\mathrm{EOS}$。
  第三步的结束候选分别为$AA\mathrm{EOS}:0.099$和
  $AB\mathrm{EOS}:0.154$；其余未结束扩展因预算耗尽不能交付。
  最佳完成文本为$B$，概率$0.32$。贪心的遗漏发生在第一步删掉$B$，
  不是第三步不会选择EOS。
\end{example}

该例中束搜索确实找到更高概率的完成序列，但$B$是否是好译文或忠实摘要，还没有任何
任务证据。模型偏好的短句可能漏条件，大束宽也可能更充分地暴露模型的这种偏好。
空文本虽然在该教学分布中合法，真实任务通常还有非空等硬条件；这些条件须进入结束检查，
不能在失败后临时改变规则来夸大搜索效果。

贪心和束搜索都没有新增可学习参数。模型概率来自既有$\bm\theta$，束宽与长度指数是搜索设置；把束宽调大不会改变已经学到的概率。束搜索适合在译文、转写或格式受限输出中保留若干高分延续，但它优化的是所声明的序列评分。开放续写中，高概率常见表达可能反复出现。Holtzman等在固定GPT-2模型与开放文本提示下观察到，最大化式生成会产生重复、过于常见的续写，低困惑度也不代表更接近人类文本的多样性。这个观察不否定翻译中的搜索用途，而是说明搜索目标须与任务质量对应；该论文的截断采样对照将在\ref{sec:nlp-generation-sampling}小节继续讨论。

\subsection{温度调整与随机选择}
\label{sec:nlp-inference-temperature}

生成指南可以有多种合法措辞。为了保留这些延续，可以在当前条件分布上执行分类抽样：生成一个$[0,1)$上的均匀数，按各词元的累计概率确定它落入的区间。每增加一个词元，条件前缀变化，下一步就要重新计算分布；这不是在固定词表权重上重复抽取。

如果希望控制当前分布的集中程度，可以先用\term{温度}$\tau>0$调整模型分数：
\begin{equation}
 q_\tau(v)=\frac{\exp(z_v/\tau)}{\sum_{w\in V}\exp(z_w/\tau)}
 =\frac{p_t(v)^{1/\tau}}{\sum_{w\in V}p_t(w)^{1/\tau}}.
 \label{eq:nlp-generation-temperature}
\end{equation}
这里$z_v$是当前词元的未归一化分数。$\tau<1$提高较大概率的相对份额，$\tau>1$使分布较分散，$\tau=1$保留原分布。零不能直接代入；最大值唯一时，零温极限才对应贪心选择，最大值并列时还需给出平局规则。

温度改变的是本次抽样所用分布，不是对事实可靠性的控制旋钮。较低温度可能稳定重复一个错误日期，较高温度可能增加合法表达，也可能增加不合法内容。判断质量仍要回到任务约束；即使只关心分布保持，也应明确保持的是原始$p_t$还是温度变换后的$q_\tau$。

温度采样不学习一个新语言模型，随机数也不承担监督。对确定字段的抽取，可以选择较集中的规则以减少无关措辞波动；对需要多种合法改写的候选生成，则可保留更多概率质量。两种用途的验收分别是字段正确与合法表达覆盖，不能用一个温度值同时保证。相同随机种子只有在词表顺序、概率计算和随机数消耗一致时才可能复现相同轨迹；加入草稿校正会改变随机数使用次序，等分布不要求逐字复现另一条运行轨迹。

\subsection{候选截断与重新归一化}
\label{sec:nlp-generation-sampling}

整个词表抽样可能取得很低概率的延续。\term{top-k采样}保留$q_\tau$最大的$k$个词元；\term{核采样}又称top-p采样，把词元按概率降序排列，保留累计概率首次达到阈值$p_0\in(0,1]$的最短前缀集合。若保留集合为$A$，实际抽样分布为
\begin{equation}
 q_A(v)=\frac{q_\tau(v)\mathbb{I}[v\in A]}{\sum_{w\in A}q_\tau(w)}.
 \label{eq:nlp-inference-truncation}
\end{equation}
集合必须非空，并预先固定边界平局规则。top-k的$k=|V|$以及top-p的$p_0=1$不截断；$p_0$很小则可能只剩一个词元。核采样动态改变集合大小，而不是保留$p_0|V|$项\scite{nlp-generation-holtzman2020}。

沿用例~\ref{ex:nlp-generation-beam}第一步的$(0.55,0.40,0.05)$。
在$\tau=1$时，top-k的$k=2$保留$A,B$，归一化为$(11/19,8/19)$。
top-p的$p_0=0.50,0.90,0.98$分别保留$\{A\}$、$\{A,B\}$、
$\{A,B,\mathrm{EOS}\}$。若先设$\tau=0.5$，平方后归一化得到
$(0.65054,0.34409,0.00538)$；此时$p_0=0.60$只保留$A$，
而原温度下同一阈值会保留$A,B$。所以比较时必须固定温度与截断次序。

一次采样还必须说明何时停。这里每条候选最多三次取样，抽到EOS立即结束，
未抽到则标为截断；只在展示时去掉EOS，不在取样时删掉它。
固定输入分布、四条候选、Python的\code{random.Random(7)}，每种方法分别重置种子，
按降序累计概率用一次均匀数选词。top-k取$k=2$得到
“$AAA$（截断）、$AA$、$AA$、$AA$”；top-p取$p_0=0.98$得到
“$AAA$（截断）、$AB$、$AB$、$AA$”。两组$\tau=1$，最多各用12次取样，
本次均实际用满12次。若教学硬条件为“禁止相邻$A$”，前组没有合格完成项，
后组只有两条$AB$合格，且它们是重复候选。这个小样本只验证取样机制，
不能作为核采样优于top-k的实验结论。

记录唯一完成候选数、重复片段、截断率和条件满足率，比只展示最喜欢的一条更完整。
真实比较还需固定模型、分词器、提示、随机数实现与候选预算。
翻译也能使用采样，但多样性仅限于合法译法；温度不会把错误日期变成可接受的创作。

top-k固定候选数量，在分布很尖锐时可能保留多余尾项，在分布平坦时又可能删掉大量合理质量；top-p固定累计质量，候选数随上下文变化。两者均不拟合参数，阈值只能依据所需任务与独立验证选择。Holtzman等的开放续写实验同时检查重复、文本多样性与人工判断，观察到核采样比最大化式生成更能保留多样表达，也避免了全词表随机采样常出现的不连贯尾部延续。其用途是开放文本续写，不是日期核验；这些质量结果不构成核采样在任意翻译或数学任务上最优的结论。

表~\ref{tab:nlp-inference-decoding-rules}把同一$c,\bm\theta,m$下的选择规则并列。比较生成质量时再固定候选条数$K$，并将未完成项计入；比较概率保持时，先确定目标到底是$p_t$还是变换后的分布。逐步贪心甚至不是完整序列概率最大化的精确算法，不能把它与从$p_t$抽样当成同一目标的两个实现。

\begin{table*}[htbp]
 \begin{tabularx}{\linewidth}{@{}p{24mm}p{36mm}p{34mm}X@{}}
  \toprule
  规则 & 选择对象与目标 & 额外学习或设定 & 交付分布与适用取舍\\
  \midrule
  贪心 & 当前概率最高词元 & 无训练；平局、停止约定 & 确定性路径；便于复现，但可能早期遗漏\\
  束搜索 & 至多$b$条前缀的序列评分 & 无训练；束宽、长度评分 & 近似搜索，不保持抽样分布；增加前缀计算\\
  温度采样 & $q_\tau$上的分类抽样 & 无训练；温度、种子 & 除$\tau=1$外改变$p_t$；保留多种措辞\\
  top-k & 最大$k$项内重新归一化 & 无训练；$k$、截断次序 & 固定集合大小；保留质量随上下文变\\
  top-p & 达到$p_0$质量的最短集合 & 无训练；阈值、边界平局 & 固定质量阈值；不保证事实或任务正确\\
  \bottomrule
 \end{tabularx}
 \caption{同一模型与输入下的解码规则。候选预算相同时，表中方法仍有不同目标；质量验收与分布保持须分别判断。}
 \label{tab:nlp-inference-decoding-rules}
\end{table*}

\section{候选生成与验证}
\label{sec:nlp-inference-verification}

前述规则可以直接逐步输出，也可以先构造候选再验证。这一环节有两种不同的用途：通过比较改变最终选中的结果，或者利用便宜草稿减少目标模型的串行调用。只有后一种用途中的特定校正规则，才带有保持目标抽样分布的保证。

\subsection{多候选生成与评分选择}
\label{sec:nlp-inference-reranking}

设按既定规则生成$K$条完整候选$y^{(1)},\ldots,y^{(K)}$，再用评分$s(c,y)$选择最高者。评分可以是模型概率，也可以是独立评判器给出的忠实性或任务通过分数。活动指南可以先排除遗漏报名条件的候选，再在合格项中比较简洁程度。生成分布影响候选覆盖，评分规则影响从已有候选中选谁，二者各自都可能成为失败来源。

如果一个二值检查器能够无误识别正确答案，且每条独立候选的成功概率为$a$，至少出现一条正确答案的概率是$1-(1-a)^K$。这是候选覆盖的计算，不代表现实评判器就能选中正确项；错误相关、重复候选和误判都会改变结果。即使评判完全正确，从$K$条中挑选也改变了交付结果的分布，不能称为保持单次抽样分布的加速。

评分的来源决定它能检查什么。JSON解析器、日期字段比对和程序单元测试使用预定规则，不需要训练；它们只能核验形式或指定输入。学习型结果验证器则把题目和完整解答作为输入，利用答案是否正确的标签拟合$s_{\bm\phi}(c,y)$，参数$\bm\phi$不同于生成器$\bm\theta$。过程验证器在每个步骤末端评分，需要步骤级标签或明确生成的继续求解反馈；其信号不能由最终答案标签自动冒充。提示一个现成模型“评价这份指南”属于调用已有参数，不等于已经训练了独立忠实性验证器。

Cobbe等在GSM8K中先训练解题生成器，再从训练题采样解答、按最终答案生成正确性标签，训练验证器并在测试时选取高分候选。论文图5观察到，验证选择随候选预算增长改善，但预算很大时受验证器误判限制；这说明它能用于可规范化答案的解题选择，却不是完美检查器。为了区分生成与选择，可以固定同一候选集合，比较概率排序、答案投票与验证器排序；候选中已经存在正确指南而选择失败时，增加生成参数并非唯一修复方向\scite{nlp-intro-cobbe2021}。

\subsection{草稿提出与精确采样校正}
\label{sec:nlp-inference-speculative}

另一种目标是交付与目标模型相同的抽样分布，同时减少它逐词调用的次数。可以让较便宜的草稿模型先提出一段词元，再让目标模型一次计算该段各位置在对应前缀下的分布。\term{推测解码}中的精确采样校正，就是逐位置决定草稿能否接受，并在拒绝时补回目标分布缺少的概率质量\scite{nlp-llm-leviathan2023,nlp-llm-chen2023speculative}。

先固定一个前缀，记目标抽样分布为$p$、草稿分布为$q$。从$q$抽取$v$后，以
\begin{equation}
 a(v)=\min\{1,p(v)/q(v)\}
 \label{eq:nlp-inference-speculative-accept}
\end{equation}
接受它；这里只会对$q(v)>0$的已抽中项求比值。接受事件给$v$的无条件概率质量为$q(v)a(v)=\min\{q(v),p(v)\}$。设拒绝概率为$D=1-\sum_w\min\{q(w),p(w)\}$。若$D>0$，拒绝后从残差分布
\begin{equation}
 r(v)=\frac{[p(v)-q(v)]_+}{\sum_w[p(w)-q(w)]_+},
 \qquad [u]_+=\max\{u,0\}
 \label{eq:nlp-inference-speculative-residual}
\end{equation}
重新抽样。两分布的概率和都为一，因此分母恰好等于$D$。最终取得$v$的概率为
\begin{equation}
 \min\{q(v),p(v)\}+D r(v)=p(v).
 \label{eq:nlp-inference-speculative-proof}
\end{equation}
若$D=0$，拒绝事件概率为零，无需计算残差分布。这也处理了$p(v)>0,q(v)=0$的词元：它们不会由草稿提出，却可以从残差中取得。

以教学分布$p=(0.6,0.3,0.1),q=(0.5,0.2,0.3)$为例，三个词元的接受概率为$1,1,1/3$，接受质量为$(0.5,0.2,0.1)$。拒绝质量为$0.2$，残差分布为$(0.5,0.5,0)$，补回后恰得$p$。拒绝后若直接从$p$再抽一次，所得质量是$(0.62,0.26,0.12)$，已经改变目标；校正不能省略。

对长度为$\gamma$的草稿，目标模型沿草稿前缀计算$\gamma$个验证分布。按位置依次执行上述接受；第一处拒绝时，先保留此前已接受的前缀，在该处按残差抽取，并丢弃其后的草稿。全部接受时，还可以从目标模型在完整草稿末端给出的分布再抽一个词元。每一步都在相同已接受前缀下保持条件分布，链式法则于是给出整段输出的分布保持。若目标规则包含温度、截断或文法约束，$p$就应是这些变换后真实要保持的分布。

这项保证依赖正确的草稿概率、目标概率及接受校正，不表示每条输出与另一随机运行逐字相同。收益还取决于草稿计算成本、目标验证成本及平均接受长度；草稿常被立即拒绝，或者验证大量无用候选时，调用轮数减少也未必缩短总耗时。

校正算法没有需要学习的参数。草稿可以来自一个独立预训练或蒸馏过的小模型，也可以来自已有统计分布；只要能给出实际提议概率，精确性就不要求$q=p$。相近的$q$主要影响接受长度，而不是校正后的正确分布。Chen等为Chinchilla训练了同词元接口、同语料的4B草稿，采用较少层数降低串行延迟；Leviathan等也试验了既有较小T5及便宜近似提议。因此“推测解码需要训练”应具体说成所选草稿的形成需要训练，不能给拒绝校正虚构损失函数。

设草稿每步成本相对普通目标步为$c_d$，验证一段的成本相对目标步为$c_v(\gamma)$，实际每轮获准交付的平均词元数为$\bar a$，忽略其他开销时，加速约为$\bar a/(\gamma c_d+c_v(\gamma))$。这只是成本分解，不是无条件实测预测。Leviathan等表2（第6页）以T5-XXL 11B为目标、T5-Small为草稿，在英译德任务上报告温度1时约$2.6$倍、贪心设置下约$3.4$倍加速。Chen等表1（第6页）在批量1、草稿长度4的Chinchilla实验中，XSum核采样的平均词元耗时由14.1毫秒降为7.52毫秒，ROUGE-2为0.112与0.114。有限随机样本的分数不必相等，等分布保证来自校正推导；两篇论文的不同目标、设备与任务也不能直接拼成快慢排名。

推测验证只问“该词元怎样按目标概率接受”，没有问“报名日期是否真实”。一个目标模型稳定偏向错误日期时，精确加速仍保留这种偏差。它适用于需要保持既有生成规则的长回答、摘要和代码生成，尤其在目标逐词计算昂贵、草稿便宜且容易被接受时；短输出和高并发计算饱和时，额外草稿与候选验证开销可能抵消收益。

\subsection{多预测头的候选构造与验证}
\label{sec:nlp-inference-multihead}

草稿不一定来自独立的小模型。Medusa在目标模型隐藏状态上接入多个预测头，各头预测不同未来位置，再把若干高分结果组成候选树，利用只允许读取同一路径祖先的注意关系验证候选\scite{nlp-llm-medusa2024}。多个头共享当前信息，却没有自动给出未来词元之间完整的条件依赖；直接把每头的最大值拼接，可能产生互不相容的续写。

第\ref{chap:nlp-large-language-models}章的多词元预测目标也可以提供未来位置的辅助头\scite{nlp-llm-mtp2024}。训练目标增加了预测监督，不意味着推断时可以无条件同时交付多个词元；仍须说明这些头怎样构造草稿，以及目标模型怎样验证。候选树的分支越多，可能覆盖更多可接受延续，也增加验证的位置数。

接受规则决定最后的性质。Medusa讨论了精确拒绝校正，也提出按目标概率阈值接受合理候选的典型接受方法。后者可以提高接受长度，但它通常改变原始抽样分布；质量相近的经验结果不能替代等分布证明。比较时须把草稿来源与接受规则分别记录。

为使未来头真正能够提出词元，Medusa把当前隐藏向量接到带残差的浅层前馈网络，再接词表投影。原语言模型头预测下一词元，第$k$个新增头预测更远的第$k+1$个词元。训练输入是已有对话文本，监督目标是实际未来词元，多个头的交叉熵按权重汇总。Medusa-1冻结基座，只学习新增头，原目标模型的概率函数不变；Medusa-2联合适配基座及预测头，并保留普通下一词元目标，采用先稳定预测头再联合更新等安排。此时目标模型本身已经变化，“相对原模型无损”不能只凭验证规则成立。

候选树将不同头的高分项组合，树注意力复用共同祖先但不读取兄弟分支。预测头监督给出未来位置的边缘信息，目标验证补上具体路径条件；两者不能互相替代。原论文在Vicuna-7B、ShareGPT训练与MT-Bench生成设置中报告，Medusa-1和Medusa-2的耗时加速分别为$2.18$倍和$2.83$倍；表2（第15页）对应质量分为基线6.17、Medusa-1为6.23、Medusa-2为6.18。该质量分来自评判模型，不是事实正确率，也不是分布相同的证明。它说明浅头能够降低这类聊天生成的串行开销，同时联合训练须检查基座质量保持。

\subsection{内部特征复用与草稿生成}
\label{sec:nlp-inference-eagle}

独立草稿模型只见词元，仍需重新推断当前上下文。EAGLE把目标模型已经得到的高层特征交给轻量草稿模块，并补入已经抽中的词元，以区分同一个历史特征之后不同抽样结果带来的分支；草稿随后经过目标模型验证\scite{nlp-llm-eagle2024}。补入已抽中词元的理由是：隐藏状态给出下一词元的分布，却未确定随机选择了哪个词元。

原始EAGLE冻结目标模型，用其高层特征序列和错开一步的词元嵌入训练浅层自回归模块。模块预测下一位置的目标特征，再通过目标模型的词表头形成草稿概率。训练同时使用预测特征与真实特征的Smooth L1误差，以及目标和草稿词元分布之间的交叉熵；特征回归提供多步回填的共同坐标，词元项使学习目标接近最终提议任务。论文还向输入特征加入小噪声，缓解连续预测误差。运行时模型生成的特征不是新观测，仍可能偏离训练时的真实特征。

原论文表2（第7页）在MT-Bench、温度1条件下，报告Vicuna-13B相对普通自回归约$2.57$倍加速，平均接受长度为3.60；同表不同模型的加速与接受长度并不严格同序，因为草稿和目标每步成本不同。该方法可以复用聊天和长解答生成中已有的目标特征，但草稿通常与具体目标检查点及隐藏维度绑定；更换目标权重以后，不能假定原有特征回归器仍然适用。

EAGLE-3改用低、中、高层特征的融合，并让草稿模块直接服务于词元预测，不再要求输出拟合目标模型的高层特征；训练时还回填自身产生的中间状态，使草稿适应连续提出多步时的输入\scite{nlp-llm-eagle32025}。这些变化影响草稿质量和成本。是否保持目标分布，仍由验证及拒绝校正规则决定，不能仅从使用了内部特征推导出来。

具体地，EAGLE-3拼接低、中、高层向量并线性压回隐藏维度，再与已选词元嵌入组合，交给草稿解码层。取消特征回归后，草稿中间向量不必等于目标特征；若训练只见真实目标向量，第二步改读自己的输出就会产生新的输入偏移。训练时测试正是将这些中间输出回填，继续计算后续词元预测目标，并以注意关系隔离不同训练分支。它不是在用户请求上更新参数，而是在离线草稿训练中模拟未来运行输入。

EAGLE-3原论文表1（第7页）在批量1、贪心、LLaMA-3.1-Instruct 8B的MT-Bench设置下，EAGLE-2为$3.16$倍，EAGLE-3为$4.40$倍；对应平均接受长度为4.05与6.13。这里固定了目标和任务，但草稿架构、目标及训练数据量共同变化，因此只能归于完整草稿方案。论文也评测推理蒸馏模型上的GSM8K生成，展示长解答加速用途；它没有因此增加数学正确率的理论保证。

表~\ref{tab:nlp-inference-verifier-contracts}总结验证对象。比较耗时时应固定目标检查点、输入、生成规则及输出上限；比较答案质量则固定候选或总预算。前者要求对已声明分布保持，后者允许改变选择，但必须证明改变有助于任务，而不是把两种收益混写。

\begin{table*}[htbp]
 \begin{tabularx}{\linewidth}{@{}p{25mm}p{42mm}p{31mm}X@{}}
  \toprule
  方法 & 参数与提议信息来源 & 验收对象 & 保证与代价\\
  \midrule
  完整候选重排 & 固定生成器；规则或另训评分器 & 答案、事实或偏好 & 改变交付分布；受覆盖和误判限制\\
  独立草稿校正 & 小模型预训练或蒸馏；校正无训练 & 目标条件词元概率 & 精确校正保持指定分布；增加草稿成本\\
  Medusa-1 & 冻结基座，未来头监督学习 & 候选路径及接受规则 & 精确规则与典型接受分开；头增加计算\\
  Medusa-2 & 基座与未来头联合适配 & 更新后目标的路径 & 须另查原质量保持；不能默认原目标不变\\
  EAGLE & 目标特征、已选词元；特征及词元目标 & 目标概率校正 & 冻结目标；特征误差影响接受长度\\
  EAGLE-3 & 多层融合、词元目标、自身状态回填 & 目标概率校正 & 冻结目标；草稿训练与验证开销另计\\
  \bottomrule
 \end{tabularx}
 \caption{草稿与验证器的职责。概率接受不是事实检查；提示评分、规则检查与独立学习型验证器也有不同的参数来源。}
 \label{tab:nlp-inference-verifier-contracts}
\end{table*}

\section{停止条件与输出约束}
\label{sec:nlp-inference-termination}

候选获得接受之后，还要确定是否继续、是否允许当前延续，以及最终结果能否交付。约束会在整个生成循环中生效，本节把它们集中说明，以免把“已经停止”误写成“已经完成”。

\subsection{结束标记与预算终止}
\label{sec:nlp-inference-stop}

模型生成EOS表示在当前分布中选择了结束词元。应用还可以设置结束字符串、最大词元数或时间上限。各停止规则的触发位置及是否保留边界标记都须明确；结束字符串可能跨越多个词元，不能只比较最后一个词元的解码文本。由于长度限制而停下的半份指南，与模型主动结束的完整指南有不同的状态。

硬性最短长度可以暂时禁止EOS，但它改变了允许延续集合，也可能迫使一个已经答完的模型继续重复。最大长度限制防止无限延续，却不保证在限制内出现完整答案。记录应保存实际生成长度和终止原因，使后续评价能够区分答错、截断和未生成。

\subsection{合法延续与结构约束}
\label{sec:nlp-inference-grammar}

若指南要交付为具有日期与地点字段的JSON，单靠提示仍可能漏引号或字段。可以维护当前文法解析状态$s_t$，确定哪些词元接入后仍能延续成合法结构，记允许集合为$A(s_t)$。把不允许项概率置零，再在允许集合内归一化，就把结构要求接入每一步选择\scite{nlp-llm-guided2023}。词元可能跨过多个字符及文法边界，因此必须检查词元对应的完整字符序列，不能把“允许下一个字符”直接当成“允许这个词元”。

生成一个词元后，解析器更新到$s_{t+1}$。EOS只有在结构已经完成时才可允许；如果$A(s_t)$为空，应报告约束或输入不相容，而不是对零质量分布抽样。允许集合还应排除无法再达到完整结构的死路，但长度预算不足仍可能使合法前缀停在未完成状态。

文法只限制形式。\code{\{"date":"2025-04-16"\}}可以是合法JSON，却可能把报名截止日期填进活动日期。语义值、事实绑定和输出用途仍需核验；文法约束也改变了抽样分布，精确推测采样应针对该受约束分布执行校正。

解析器的状态转移来自指定文法和词元到字符的映射，不需要额外训练。引导生成研究将允许字符串组织为状态机，并预计算状态允许的词元，避免每步重新逐项试解析；对标点或固定字段名这样的唯一合法延续，还可以减少不必要的模型选择。这提供了结构化数据生成的具体用途，但省去的只是固定结构上的重复工作，无法决定字段值是否有证据，也不能承诺任意文法下都同样加速。

\subsection{完成状态与输出核验}
\label{sec:nlp-inference-output-check}

结束之后，先按分词器恢复文本并删除约定的控制标记，再分别核对终止原因、结构完整性和任务要求。活动指南需要包含活动时间、地点与校外报名要求；具备两个句号不保证这三项已经保留。检查器可以给出遗漏字段或冲突来源的位置，以便继续、重新生成或转入失败处理。

再次生成与选择也会消耗预算并改变最终输出的分布。若只展示多次尝试中最好的一条，却按单次调用报告成本，就把选择规则的收益归给了模型。应保留尝试次数、候选及检查依据；超过预定修订预算仍未完成时，明确返回未完成。任务正确性如何核验将在应用章展开，第\ref{chap:nlp-evaluation}章统一规定怎样比较这些运行规则。

\section{本章小结}
\label{sec:nlp-inference-summary}

固定模型参数以后，输出仍由输入条件、状态计算、选择规则和终止约束共同决定。因果缓存利用历史可见关系保持不变这一条件减少重复计算；掩码恢复则采用不同的迭代循环。搜索、温度与截断采样改变候选选择，重排序进一步改变交付分布；精确推测采样通过接受和残差校正保持指定条件分布，其加速仍受草稿与验证成本限制。

这些机制适用于写作、翻译和多步求解，但它们本身没有保证结果正确。下一章以推理能力为专题，把各训练阶段提供的知识、轨迹和反馈与推断阶段的计算结合起来，解释模型怎样形成并改进可核验的求解行为。

\paragraph{延展阅读}
退化文本与核采样的分析适合继续研究候选分布怎样受截断规则影响；推测采样的两项工作
给出接受与残差校正，EAGLE则从草稿特征预测扩展加速路线。比较这些方法时，应同时报告
目标分布是否保持、验证模型调用次数、候选接受率和端到端耗时，不能只用每轮产生的
词元数判断推断收益
\scite{nlp-generation-holtzman2020,nlp-llm-leviathan2023,nlp-llm-chen2023speculative,nlp-llm-eagle2024}
